\section{What Was Defined in 1336, and What It Cost}

The best way into the Catholic doctrine of heaven is not through a catechism paragraph but through a quarrel, because the quarrel shows exactly where the load-bearing walls are.

\subsection{A pope preaching the other side}

On All Saints' Day of 1331, and in further sermons over the following two years, Pope John XXII taught at Avignon that the souls of the just do not see the divine essence immediately after death. On his account they wait---consoled, in the presence of Christ's humanity, but not yet admitted to the vision of God as he is---until the resurrection of the body and the general judgment. He had held the opinion before his election and continued to hold it as pope.\endnote{\emph{The Catholic Encyclopedia}, art.\ ``Pope John XXII'' (vol.~8, New York 1910), read 2026-07-26 at New Advent: ``Before his elevation to the Holy See, he had written a work on this question, in which he stated that the souls of the blessed departed do not see God until after the Last Judgment. After becoming pope, he advanced the same teaching in his sermons.'' The article is a public-domain 1910 reference work used here as a historical witness for the sequence of events, not as a doctrinal authority; the acts themselves (John's \emph{Ne super his} and Benedict XII's \emph{Benedictus Deus}) are quoted below from the received Latin.}

The reaction was severe and instructive. Theologians called the view heretical. The University of Paris debated it; in December 1333 its theologians decided in favour of the immediate vision and petitioned the pope to confirm their decision, while pointing out---this is the part worth noticing---that the pope had given no decision on the question but had only advanced his personal opinion.\endnote{\emph{Catholic Encyclopedia}, art.\ ``Pope John XXII,'' read 2026-07-26: John ``wrote to King Philip IV on the matter (November, 1333), and emphasized the fact that, as long as the Holy See had not given a decision, the theologians enjoyed perfect freedom in this matter''; the Paris theologians ``pointed out that the pope had given no decision on this question but only advanced his personal opinion.'' In a consistory of 3 January 1334 John ``explicitly declared that he had never meant to teach aught contrary to Holy Scripture or the rule of faith and in fact had not intended to give any decision whatever.'' The distinction being drawn---between what a pope says and what a pope defines---is the ordinary Catholic one and is not this article's invention; nothing here is a judgment about the man.} John himself had written to the King of France in November 1333 insisting that so long as the Holy See had given no decision, theologians were perfectly free on the matter.

On 3 December 1334, the day before he died, John XXII issued the bull \emph{Ne super his}. It is a short and rather careful document, and its central sentence deserves to be read slowly:

\begin{studyframe}
\latin{Fatemur siquidem et credimus, quod animae purgatae separatae a corporibus sunt in caelo, caelorum regno et paradiso et cum Christo in consortio angelorum congregatae et vident Deum de communi lege ac divinam essentiam facie ad faciem clare, in quantum status et condicio compatitur animae separatae.}\endnote{John XXII, bull \emph{Ne super his}, 3 December 1334, as received at Denzinger nn.~990--991 (older numbering 457a--457b); Latin read 2026-07-26 at the patristica.net web presentation of the Latin \emph{Enchiridion symbolorum}, the edition registered in the source library for this witness. Working rendering (this article's, governed by the Latin, and labelled as a gloss, not as an approved translation): ``We confess and believe that souls, purified and separated from their bodies, are gathered in heaven, in the kingdom of heaven and paradise, and with Christ in the company of the angels, and see God and the divine essence face to face clearly, according to the common law, \emph{so far as the state and condition of a separated soul allows}.'' The whole document then subjects everything he has said or written to the determination of the Church and of his successors.}
\end{studyframe}

He confesses that the purified souls do see God and the divine essence face to face, clearly. And then he adds seven words that a lawyer would circle: \latin{in quantum status et condicio compatitur animae separatae}---so far as the state and condition of a separated soul allows. Whether that is a genuine retraction or a retraction with the door left ajar is a question this article does not need to settle. What matters is that everyone at the time could see the door.

\subsection{The definition}

His successor closed it. Jacques Fournier---a Cistercian, a Paris doctor of theology, and, by the account of the standard older reference, a man who heard out the defenders of the deferred vision and then spent four months with a commission of theologians in patristic research before publishing---was elected on 20 December 1334 and promulgated the apostolic constitution \emph{Benedictus Deus} on 29 January 1336.\endnote{\emph{The Catholic Encyclopedia}, art.\ ``Benedict XII'' (vol.~2, New York 1907), read 2026-07-26 at New Advent: ``Eager to solve the question, Benedict heard the opinions of those maintaining the theory of deferred vision, and, with a commission of theologians, gave four months to patristic research. Their labours terminated in the proclamation (29 January, 1336) of the Bull `Benedictus Deus' defining the immediate intuitive vision of God by the souls of the just having no faults to expiate.'' The dates of election and enthronement are from the same article.} The doctrinal core, as it is received in the standard collection, runs:

\begin{studyframe}
\latin{Hac in perpetuum valitura Constitutione auctoritate Apostolica diffinimus: quod secundum communem Dei ordinationem animae sanctorum omnium\ldots\ etiam ante resumptionem suorum corporum et iudicium generale\ldots\ fuerunt, sunt et erunt in caelo, caelorum regno et paradiso caelesti cum Christo, sanctorum Angelorum consortio congregatae, ac post Domini Jesu Christi passionem et mortem viderunt et vident divinam essentiam visione intuitiva et etiam faciali, nulla mediante creatura in ratione obiecti visi se habente, sed divina essentia immediate se nude, clare et aperte eis ostendente, quodque sic videntes eadem divina essentia perfruuntur, necnon quod ex tali visione et fruitione eorum animae, qui iam decesserunt, sunt vere beatae et habent vitam et requiem aeternam.}\endnote{Benedict XII, apostolic constitution \emph{Benedictus Deus}, 29 January 1336, as received at Denzinger nn.~1000--1001 (older numbering 530); Latin read 2026-07-26 at the registered patristica.net presentation of the Latin \emph{Enchiridion}. The ellipses omit the constitution's long enumeration of the classes of soul covered (the saints of before the Passion, apostles, martyrs, confessors, virgins, the baptised faithful in whom nothing was to be purged, those purged after death, and baptised children dying before the use of free will) and its clause about the Ascension. The full enumeration is summarised in the prose above; nothing doctrinally material is omitted. The 1336 Latin act is public domain; printed editions of Denzinger retain their own editorial rights, and the web presentation used here does not identify the exact printing it transcribes, which is recorded as this article's verification ceiling for the wording.}
\end{studyframe}

The official English of the central clause, as the Catechism quotes it, reads: these souls ``have seen and do see the divine essence with an intuitive vision, and even face to face, without the mediation of any creature.''\endnote{\emph{Catechism of the Catholic Church} 1023, quoting \emph{Benedictus Deus} (DS 1000), official Vatican English web text, read 2026-07-26. The Catechism's excerpt begins ``By virtue of our apostolic authority, we define the following'' and preserves the constitution's enumeration in abbreviated form with ellipses; the wording quoted here is the Catechism's own rendering and is not altered.}

Read the Latin phrase by phrase and the definition turns out to be doing several distinct jobs at once.

\emph{\latin{Secundum communem Dei ordinationem}}---according to God's common ordering. The constitution defines what ordinarily happens, and does not pronounce on whether God may dispose otherwise in a particular case. It is a general law, stated as a general law.

\emph{\latin{Etiam ante resumptionem suorum corporum et iudicium generale}}---even before they take up their bodies again and before the general judgment. This is the point of the whole document, the exact proposition John XXII had denied. The soul does not wait for the resurrection to see God.

\emph{\latin{Post Domini Jesu Christi passionem et mortem}}---since the Passion and death of the Lord. The vision is not a natural entitlement of separated souls; it dates from a historical event. The just of the old covenant, expressly included in the enumeration, saw nothing of the kind before Calvary.

\emph{\latin{Viderunt et vident divinam essentiam}}---have seen and do see the divine essence. Not God's effects; not God's glory as distinct from God; the divine essence.

\emph{\latin{Visione intuitiva et etiam faciali}}---by an intuitive vision, and even facial. ``Intuitive'' excludes the discursive, inferential knowing by which a mind reasons from effects to a cause. ``Facial'' is Saint Paul's word, taken from ``we see now through a glass in a dark manner; but then face to face.''\endnote{1 Cor 13:12, tracked verse text of the Douay--Rheims, read 2026-07-26. The Latin \latin{facialis} in the constitution is the scholastic adjective formed from the Vulgate's \latin{facie ad faciem}; it is a technical term of art, not an assertion that God has a face, a point Augustine had already had to make explicitly (\emph{De civitate Dei} XXII.29: ``By `the face' of God we are to understand His manifestation, and not a part of the body similar to that which in our bodies we call by that name'').}

\emph{\latin{Nulla mediante creatura in ratione obiecti visi se habente}}---no creature intervening \emph{in the role of an object seen}. This is the most carefully machined phrase in the document, and the qualifier is everything. What is excluded is a created \emph{object} standing between the mind and God, some image or likeness or delegated glory that the blessed would see instead of God. What is not excluded, because it is not the same thing, is a created \emph{capacity} in the seer by which God himself is seen. That distinction is not a later evasion; it is the drafting, and the next section shows a theologian occupying the space it leaves open sixty years before the definition was written.

\emph{\latin{Sunt vere beatae et habent vitam et requiem aeternam}}---they are truly blessed and have eternal life and rest. The constitution goes on to add that this vision and enjoyment ``void the acts of faith and hope in them, inasmuch as faith and hope are properly theological virtues,'' and that once begun it continues without any interruption up to the last judgment and thereafter for ever.\endnote{\emph{Benedictus Deus}, DS 1001: \latin{ac quod visio huiusmodi divinae essentiae eiusque fruitio actus fidei et spei in eis evacuant, prout fides et spes propri(a)e theologicae sunt virtutes; quodque, postquam inchoata fuerit vel erit talis intuitiva ac facialis visio et fruitio in eisdem, eadem visio et fruitio sine aliqua intercisione (intermissione) seu evacuatione praedictae visionis et fruitionis continuata exstitit et continuabitur usque ad finale iudicium et ex tunc usque in sempiternum.} The parenthesised variants are the transcription's own and are reproduced rather than silently resolved. The final paragraph of the constitution (DS 1002) defines the immediate descent to hell of those who die in actual mortal sin and the general judgment in which all will appear with their bodies; that subject belongs to the other volumes of this series and is not developed here.}

\subsection{What the definition does not say}

An honest reading has to be as attentive to the silences as to the words, and here the silences are enormous.

The constitution says nothing about how a finite mind performs an act of which it is not naturally capable. It says nothing about what the vision is \emph{like}. It says nothing about whether the vision admits of degrees. It says nothing about what the risen body contributes, beyond noting that the vision precedes the body's return. It says nothing about place: heaven appears as \latin{caelum}, \latin{caelorum regnum}, \latin{paradisus caelestis}, and those are the terms of the received formula, not a topography. It says nothing about the company of the blessed except that they are gathered ``in the fellowship of the holy angels''---a phrase inherited from John XXII's own bull two years earlier. It says nothing about time, activity, memory, recognition, growth, or any of the questions a bereaved person actually asks.

This is worth dwelling on because it establishes the proportions of everything that follows. On the single most important question about heaven the Church has spoken with maximum solemnity and minimum extent. Everything else in the received picture---and it is a very great deal---either comes from Scripture read as Scripture, or from the reasoning of theologians, or from nowhere.

\subsection{Two more fixed points}

Two later definitions bear directly on this article's subject, and both should be on the table from the start.

The Council of Florence, in the bull of union with the Greeks \emph{Laetentur caeli} of 6 July 1439, defined that the souls of those who die with no stain of sin, and those purified after death, ``are received at once into heaven and see clearly God himself, three and one, as he is---yet in accordance with the diversity of merits, one more perfectly than another.''\endnote{Council of Florence, bull \emph{Laetentur caeli}, 6 July 1439, DS 1305 (older numbering 693): \latin{in caelum mox recipi et intueri clare ipsum Deum trinum et unum, sicuti est, pro meritorum tamen diversitate alium alio perfectius}. Latin read 2026-07-26 at the registered patristica.net presentation; the English is this article's labelled working gloss governed by the Latin. Note that \emph{Laetentur caeli} is a decree of union with the Greek Church, subscribed at Florence and afterwards repudiated at Constantinople; its doctrinal force in the Catholic Church does not depend on its reception in the East, and this article does not treat the repudiation, which lies outside its scope.} Two additions to 1336 are made there, and both are consequential: the vision is of God \latin{sicuti est}, ``as he is,'' the phrase of the First Epistle of John;\endnote{1 Jn 3:2: ``We know that when he shall appear we shall be like to him: because we shall see him as he is.'' Tracked verse text, read 2026-07-26.} and it is unequal---\latin{alium alio perfectius}, one more perfectly than another, in proportion to merits. The degrees of glory are not a pious embellishment. They are defined.

The Fourth Lateran Council, in the profession of faith \emph{Firmiter} of 1215, defined of the resurrection only this: all ``will rise again with their own bodies which they now bear, to receive according to their works, whether they have been good or evil.''\endnote{Fourth Lateran Council (1215), constitution 1, \emph{Firmiter credimus}, DS 801 (older numbering 429): \latin{qui omnes cum suis propriis resurgent corporibus, quae nunc gestant, ut recipiant secundum opera sua, sive bona fuerint sive mala, illi cum diabolo poenam perpetuam, et isti cum Christo gloriam sempiternam.} Latin read 2026-07-26 at the registered patristica.net presentation; English gloss labelled. The Catechism quotes the same clause at CCC 999, in the form ``all of them will rise again with their own bodies which they now bear.''} \latin{Cum suis propriis corporibus, quae nunc gestant}: with their own proper bodies, the ones they are carrying around at present. That is the whole of it. Every further specification---the qualities of the risen body, its age, its powers, its appearance---is theology, and the section on the risen body below will show how much theology there is and how frankly its best authors admitted what they were doing.

\subsection{What a definition is for}

It is tempting to read the 1336 affair as an embarrassment to be explained away, and equally tempting to read it as a scandal. It is neither; it is a demonstration of how the machinery is supposed to work and, on this occasion, did. A pope held and preached an opinion. Theologians said he was wrong, in public, at length, and he conceded that until the Holy See had decided they were free to say so. He was not asked to be infallible about something he had not defined, and he did not claim to be. When a definition finally came, it came from a successor, after a commission and four months in the Fathers, and it was drafted with the precision of a legal instrument by a man who had studied the objections.

The lesson for a reader of this article is the operative one. The Catholic Church has said less about heaven than her devotional literature suggests, and what she has said she has said in language that repays being read as carefully as it was written. The rest of this study is an attempt to take the rest of the tradition just as seriously---which means, above all, refusing to give it the weight of the paragraphs just quoted.
